\documentclass[landscape,a0paper,fontscale=0.292]{baposter}

\usepackage[vlined]{algorithm2e}
\usepackage{times}
\usepackage{calc}
\usepackage{url}
\usepackage{graphicx}
\usepackage{amsmath}
\usepackage{amssymb}
\usepackage{relsize}
\usepackage{multirow}
\usepackage{booktabs}

\usepackage{graphicx}
\usepackage{multicol}
\usepackage[T1]{fontenc}
\usepackage{ae}
\usepackage{enumitem}

\usepackage{colortbl}
\usepackage{xcolor}
\graphicspath{{images/}}

\setlist[itemize]{leftmargin=*,nosep}
 \setlength{\columnsep}{0.7em}
 \setlength{\columnseprule}{0mm}


% %%%%%%%%%%%%%%%%%%%%%%%%%%%%%%%%%%%%%%%%%%%%%%%%%%%%%%%%%%%%%%%%%%%%%%%%%%%%%%%%
% % Save space in lists. Use this after the opening of the list
% %%%%%%%%%%%%%%%%%%%%%%%%%%%%%%%%%%%%%%%%%%%%%%%%%%%%%%%%%%%%%%%%%%%%%%%%%%%%%%%%
 \newcommand{\compresslist}{%
 \setlength{\itemsep}{0pt}%
 \setlength{\parskip}{0pt}%
 \setlength{\parsep}{0pt}%
 }
\renewcommand{\rmdefault}{ptm} % Arial
\renewcommand{\sfdefault}{ptm} % Arial

%%%%%%%%%%%%%%%%%%%%%%%%%%%%%%%%%%%%%%%%%%%%%%%%%%%%%%%%%%%%%%%%%%%%%%%%%%%%%
%% Begin of Document
%%%%%%%%%%%%%%%%%%%%%%%%%%%%%%%%%%%%%%%%%%%%%%%%%%%%%%%%%%%%%%%%%%%%%%%%%%%%%
\begin{document}
%%%%%%%%%%%%%%%%%%%%%%%%%%%%%%%%%%%%%%%%%%%%%%%%%%%%%%%%%%%%%%%%%%%%%%%%%%%%%
%% Here starts the poster
%%---------------------------------------------------------------------------
%% Format it to your taste with the options
%%%%%%%%%%%%%%%%%%%%%%%%%%%%%%%%%%%%%%%%%%%%%%%%%%%%%%%%%%%%%%%%%%%%%%%%%%%%%
\begin{poster}{
 % Show grid to help with alignment
 grid=false,
 columns=5,
 % Column spacing
 colspacing=0.7em,
 % Color style
 headerColorOne=cyan!20!white!90!black,
 borderColor=cyan!30!white!90!black,
 % Format of textbox
 textborder=faded,
 % Format of text header
 headerborder=open,
 headershape=roundedright,
 headershade=plain,
 background=none,
 bgColorOne=cyan!10!white,
 headerheight=0.12\textheight}
 % Eye Catcher
 {
      \includegraphics[width=0.08\linewidth]{PolyU_logo.png}
      \makebox[0.01\textwidth]{} 
      \raisebox{0.08\height}{\includegraphics[width=0.10\linewidth]{CUHK_logo.jpg}}
      \makebox[0.04\textwidth]{} 
 }
 % Title
 {\sc\Huge\bf A Trilateral Weighted Sparse Coding Scheme \\ for Real-World Image Denoising}
 % Authors
 {
 \vspace{0.1em} 
 \normalsize{Jun Xu$^1$, Lei Zhang$^1$, David Zhang$^{1,2}$ }
 \\[0.1em]
 {\small
 $^1$ Department of Computing, Hong Kong Polytechnic University, Hong Kong SAR, China
 \\
 $^2$ School of Science and Engineering, The Chinese University of Hong Kong (Shenzhen), Shenzhen, China
 \\[0.1em] 
 }
 }
 % University logo
 {
    \begin{tabular}{r}
        \includegraphics[width=0.15\linewidth]{cropped-Logo_ECCV18_RGB_1200x280}
    \end{tabular}
 }
%%%%%%%%%%%%%%%%%%%%%%%%%%%%%%%%%%%%%%%%%%%%%%%%%%%%%%%%%%%%%%%%%%%%%%%%%%%%%%
%%% Now define the boxes that make up the poster
%%%---------------------------------------------------------------------------
%%% Each box has a name and can be placed absolutely or relatively.
%%% The only inconvenience is that you can only specify a relative position 
%%% towards an already declared box. So if you have a box attached to the 
%%% bottom, one to the top and a third one which should be inbetween, you 
%%% have to specify the top and bottom boxes before you specify the middle 
%%% box.
%%%%%%%%%%%%%%%%%%%%%%%%%%%%%%%%%%%%%%%%%%%%%%%%%%%%%%%%%%%%%%%%%%%%%%%%%%%%%%

%%%%%%%%%%%%%%%%%%%%%%%%%%%%%%%%%%%%%%%%%%%%%%%%%%%%%%%%%%%%%%%%%%%%%%%%%%%%%%
\headerbox{\bf\color{blue} Problem Definition and Contribution}{name=contribution,column=0,row=0,span=2}{
    \textbf{\color{blue}Goal:} Estimating the surface normal of a static object given multiple images captured under varying light directions.
    \begin{center}
        \vspace{-0.8em}
        \centering\includegraphics[width=0.8\linewidth]{images/intro}
    \end{center}

    \vspace{-0.8em}
    \textbf{\color{blue}Key Contributions:}
    A flexible deep learning framework for photometric stereo that 
    \begin{itemize}
        \item does not depend on a pre-defined set of light directions during training and testing.
        \item processes an arbitrary number of input images in an order-agnostic manner.
        \item generalizes well on real data after training only on the synthetic data.
        \item achieves state-of-the-art results in calibrated photoemtric stereo and promising results in uncalibrated scenario.
    \end{itemize}  
}

%%%%%%%%%%%%%%%%%%%%%%%%%%%%%%%%%%%%%%%%%%%%%%%%%%%%%%%%%%%%%%%%%%%%%%%%%%%%%%
\headerbox{\bf\color{blue} Experiments \& Results}{name=results,column=2,row=0,span=3}{
    \begin{minipage}[t]{0.5\textwidth}
        \textbf{\color{blue}Synthetic Datasets for Training:} 
        \vspace{-0.2em}
        \begin{center}
            \includegraphics[width=\textwidth]{images/datasets.pdf}
        \end{center}
    \end{minipage}
    \begin{minipage}[t]{0.5\textwidth}
        \textbf{\color{blue}Quantitative Results on DiLiGenT Main Dataset:} 
        \vspace{-0.2em}
        \begin{center}
            \includegraphics[width=0.98\textwidth]{images/res_quant_diligent_main}
        \end{center}
    \end{minipage}

    \vspace{0.5em}

    \begin{minipage}[t]{0.525\textwidth}
        \textbf{\color{blue}Feature Visualization:}
        \vspace{-0.5em}
        \begin{center}
            \includegraphics[width=\textwidth]{images/visualization}
        \end{center}
        \vspace{-0.9em}
        \begin{itemize}
            \item Different regions with similar normal directions are fired in different channels. Each channel can therefore be interpreted as the probability of the normal belonging to a certain direction.
        \end{itemize}
    \end{minipage}
    \hfill
    \begin{minipage}[t]{0.465\textwidth}
    \textbf{\color{blue}Qualitative Results on DiLiGenT Main Dataset:} 
        \begin{center}
            \vspace{-0.5em}
            \includegraphics[width=\textwidth]{images/res_qual_diligent_main}
        \end{center}
    \end{minipage}

    \vspace{0.8em}
    \textbf{\color{blue}Qualitative Results on the Gourd\&Apple Dataset and Light Stage Data Gallery:}
    \vspace{-0.8em}
    \begin{center}
        \includegraphics[width=1\textwidth]{images/gourd_stage}
    \end{center}

    \begin{minipage}[t]{0.6\textwidth}
        \textbf{\color{blue}Quantitative Results on Spheres Rendered with 100 Different Materials:} 
        \vspace{-0.5em}
        \begin{center}
            \includegraphics[width=\textwidth]{images/100brdf.png}
        \end{center}
    \end{minipage}
    \begin{minipage}[t]{0.4\textwidth}
        \textbf{\color{blue}Quantitative Results of Uncalibrated PS-FCN on DiLiGenT Main Dataset:} 
        \vspace{-0.5em}
        \begin{center}
            \includegraphics[width=\textwidth]{images/res_uncalibrated}
        \end{center}
        \vspace{-1.5em}

        \begin{center}
        \begin{minipage}{0.56\linewidth}
            \begin{center}
            \textbf{Project Webpage}: \\
            \vspace{0.5em}\textbf{Code} \& \textbf{Dataset} \& \textbf{Model}
            \end{center}
        \end{minipage}
        \begin{minipage}{0.24\linewidth}
            \begin{center}
                \includegraphics[width=\linewidth]{images/PS-FCN_QRCode.png}
            \end{center}
        \end{minipage}
        \end{center}
    \end{minipage}
}

\headerbox{\bf\color{blue} Formulation}{name=formulation,column=0,below=contribution,span=2}{
    \textbf{\color{blue}Assumption:} Orthographic projection and directional lights.
    
    \textbf{\color{blue}Image Formation Equation:} Given $q$ color images of an object with $p$ pixels captured under different light directions, a normal matrix $\mathbf{N}_{3\times p}$, a light direction matrix $\mathbf{L}_{3\times q}$, and an observation matrix $\mathbf{I}_{3\times p\times q}$ can be constructed. Denoting the BRDFs for all observations as $\boldsymbol{\Theta}_{3\times p\times q}$, the image formation equation can be written as
    \vspace{-0.5em}
    \begin{equation}
        \textbf{I} = \boldsymbol{\Theta} \circ \text{repmat}(\textbf{N}^\top \textbf{L}, 3),
        \vspace{-0.5em}
    \end{equation}
    where $\circ$ represents element-wise multiplication, and $\text{repmat}(\mathbf{X},3)$ repeats the matrix $\mathbf{X}$ three times along the first dimension. 
    
    \textbf{\color{blue}Main Idea:} Our method directly learns the mapping from $(\textbf{I}, \textbf{L})$ to $\textbf{N}$ without explicitly modeling $\boldsymbol\Theta$.
}

%%%%%%%%%%%%%%%%%%%%%%%%%%%%%%%%%%%%%%%%%%%%%%%%%%%%%%%%%%%%%%%%%%%%%%%%%%%%%%
\headerbox{\bf\color{blue} Method}{name=abstract,column=0,below=formulation,span=2}{
    \textbf{\color{blue}Network Architecture:} PS-FCN consists of three components, namely a shared-weight feature extractor, a fusion layer, and a normal regression network.
    \vspace{-0.5em}
    \begin{center}
     \includegraphics[width=0.8\textwidth]{images/network_v2.pdf}
    \end{center}
    \vspace{-0.5em}
    \textbf{\color{blue}Loss function:}
    \vspace{-0.5em} 
    \begin{equation}
        L_{normal} = \frac{1}{hw} \sum_{i,j} (1 - \mathbf{N}_{ij} \cdot \tilde{\mathbf{N}}_{ij})
        \vspace{-0.5em} 
    \end{equation}
    where $\mathbf{N}_{ij}$ and $\tilde{\mathbf{N}}_{ij}$ denote the predicted normal and the ground truth, respectively.
}

\end{poster}
\end{document}
